% f03_claim1 variant b -- three fabrics, one root box.
% Data: tables/invariance.tex (sixteen leaf counts, one 64-bit root each,
%       divergent column all zero; the N=128 row supplies 0x415F7E8FF9581062
%       and its 371 passing HW configurations, and the N=512 row 456).
%       \WallTBs{} (data/coverage_macros.tex), \InvNonSettling{},
%       \InvLeafCounts{}, \InvMultiPoint{} (data/results_macros.tex),
%       \RouteTopoCount{}, \RouteRuns{}, \RouteReexecInst{}, \RouteBitExact{}
%       (data/wround_macros.tex, the re-execution reported in
%       sec:results-invariance).
% Strategy: CONVERGENCE. Where a keeps three root boxes side by side and lets
% the reader compare them, b draws physically one box that all three fabrics
% write into, so the claim reads as a single arrow head rather than a triple
% coincidence -- and the campaign counts that back it sit in that one place.
% Colour: the fabric is incidental machinery (spSky nodes, spGrid links); the
% route -- the thing that differs between panels -- is spPurple, kept one hue so
% that dash pattern remains the panel-distinguishing channel exactly as the
% captions say; the single root box every panel writes into is the confirmed
% invariant (spGreen).
% Target: IEEE double column, 7.16in = 18.2cm. Drawn inside 15.4cm.
\begin{tikzpicture}[
  x=1cm, y=1cm,
  pe/.style={draw=spSky!80!spInk, fill=spSkyF, circle, line width=0.3pt,
             minimum size=0.28cm, inner sep=0pt},
  sw/.style={draw=spSky!80!spInk, fill=spSky!35, rectangle, line width=0.3pt,
             minimum width=0.32cm, minimum height=0.20cm, inner sep=0pt},
  lk/.style={draw=spGrid, line width=0.35pt},
  rte/.style={draw=spPurple!85!spInk, line width=1.0pt},
  rteb/.style={draw=spPurple!85!spInk, line width=1.2pt, dash pattern=on 2pt off 1.2pt},
  rtec/.style={draw=spPurple!85!spInk, line width=1.0pt, dotted},
  ar/.style={draw=spGreen, line width=0.5pt, -{Latex[length=1.8mm,width=1.2mm]}},
  lbl/.style={font=\tiny, align=center, inner sep=1pt, text=spInk},
  pan/.style={draw=spGrid, line width=0.5pt, rounded corners=1.5pt, fill=spGrid!45},
]
% ================= panel 1: ring ==========================================
\draw[pan] (0.30,1.75) rectangle (4.60,4.35);
\node[lbl, anchor=north] at (2.45,4.28) {\textbf{ring}, $G{=}8$};
\foreach \i in {0,...,7}{
  \pgfmathsetmacro{\a}{90-\i*45}
  \node[pe] (R\i) at ({2.45+0.95*cos(\a)},{2.90+0.62*sin(\a)}) {};
}
\foreach \i/\j in {0/1,1/2,2/3,3/4,4/5,5/6,6/7,7/0}{ \draw[lk] (R\i) -- (R\j); }
\draw[rte] (R4) -- (R5) -- (R6) -- (R7) -- (R0);
\node[lbl, anchor=south, text=spPurple!85!spInk] at (2.45,1.80) {solid route: 4 hops};

% ================= panel 2: torus =========================================
\draw[pan] (5.35,1.75) rectangle (9.65,4.35);
\node[lbl, anchor=north] at (7.50,4.28) {\textbf{2D torus}, $G{=}9$};
\foreach \i in {0,1,2}{ \foreach \j in {0,1,2}{
  \node[pe] (T\i\j) at ({6.80+\i*0.70},{2.25+\j*0.62}) {};
}}
\foreach \j in {0,1,2}{ \draw[lk] (T0\j) -- (T1\j) -- (T2\j); }
\foreach \i in {0,1,2}{ \draw[lk] (T\i0) -- (T\i1) -- (T\i2); }
\foreach \j in {0,1,2}{ \draw[lk] (T0\j) to[out=170,in=10,looseness=1.25] (T2\j); }
\foreach \i in {0,1,2}{ \draw[lk] (T\i0) to[out=260,in=100,looseness=1.25] (T\i2); }
\draw[rteb] (T00) -- (T10) -- (T20) -- (T21) -- (T22);
\node[lbl, anchor=south, text=spPurple!85!spInk] at (7.50,1.80) {dashed route: dimension-ordered};

% ================= panel 3: fat-tree ======================================
\draw[pan] (10.40,1.75) rectangle (15.10,4.35);
\node[lbl, anchor=north] at (12.75,4.28) {\textbf{fat-tree}, $G{=}8$};
\node[sw] (FR) at (12.75,3.55) {};
\node[sw] (FA) at (11.85,3.05) {};
\node[sw] (FB) at (13.65,3.05) {};
\draw[lk] (FA) -- (FR) -- (FB);
\foreach \i in {0,...,3}{
  \pgfmathsetmacro{\x}{11.30+\i*0.37}
  \node[pe] (FL\i) at (\x,2.35) {};
  \draw[lk] (FL\i) -- (FA);
}
\foreach \i in {0,...,3}{
  \pgfmathsetmacro{\x}{13.10+\i*0.37}
  \node[pe] (FM\i) at (\x,2.35) {};
  \draw[lk] (FM\i) -- (FB);
}
\draw[rtec] (FL0) -- (FA) -- (FR) -- (FB) -- (FM3);
\node[lbl, anchor=south, text=spPurple!85!spInk] at (12.75,1.80) {dotted route: up--over--down};

% ================= convergence ============================================
\draw[ar] (2.45,1.68) -- (2.45,1.30) -- (6.20,1.30) -- (6.20,1.02);
\draw[ar] (7.50,1.68) -- (7.50,1.02);
\draw[ar] (12.75,1.68) -- (12.75,1.30) -- (8.80,1.30) -- (8.80,1.02);
\node[draw=spGreen, fill=spGreenF, line width=0.8pt, inner sep=3.5pt,
      rounded corners=1.5pt, align=center, font=\tiny, text=spInk] (ROOT) at (7.50,0.66)
  {the same 64-bit root\\[1pt]\texttt{0x415F7E8FF9581062} \ ($N{=}128$)};

\node[lbl, anchor=north, text width=15.0cm, text=spInk] at (7.50,0.22)
  {One root per leaf count, not one root per route.
   \InvLeafCounts{} leaf counts, \InvLeafCounts{} roots, divergence column
   zero throughout the invariance table; \InvMultiPoint{} of the \WallTBs{}
   settled simulations lie at leaf counts carrying two or more
   $(G,\text{topology})$ points. A re-execution of every passing configuration
   at $G \le \RouteGMax$ --- \RouteRuns{} runs over \RouteReexecInst{} suite
   instances spanning all \RouteTopoCount{} topologies --- agrees on one
   distinct root per instance.};
\end{tikzpicture}
